\section{Ontology Governance and the Co-responsibility Paradigm}
\label{sec:governance}

\subsection{The Non-Delegability Premise}
\label{subsec:non_delegability}

The Co-responsibility Architecture rests on a premise that is legal before it
is technical: \textbf{responsibility for a financial record cannot be delegated
to a non-human agent, whether deterministic or stochastic.} A software
artifact cannot be a defendant. It cannot be fined, sanctioned, struck off a
professional register, or imprisoned. As the University of Waterloo Centre for
Information System Assurance frames the accountability crisis in AI assurance,
``you cannot fine an algorithm, imprison a neural network, or shame a
dataset'' (Tucker \& Kapoor, 2025). The legal accountability for every posted
transaction therefore rests, necessarily and at all times, on an identifiable
human professional.

This position is not merely our engineering preference; it converges with
warnings from both ends of the AI debate. Arguing against live legislative
proposals to grant AI agents legal personhood, Yuval Noah Harari warns that doing
so would hand them ``a master key'' to the financial, economic, and political
systems, and that ``countries that grant legal personhood to AIs risk becoming
something for which the historical record offers no analogy\dots\ a country whose
inhabitants could be governed by non-human corporations'' (Harari, \emph{We
Should Not Grant Legal Personhood to AI Agents}, Financial Times, 9~June~2026).
From the technical side, Yann LeCun observes that today's language models cannot
reliably anticipate the effects of their own actions: ``it's a very bad way to
produce an action if you're not able to predict the consequences of it. In fact,
it might be dangerous'' --- and on that basis judges agentic systems built on
pure LLMs unfit to act unsupervised (LeCun, lecture at Brown University,
1~April~2026). An entity that can be granted neither legal standing nor a reliable
model of the consequences of its own actions cannot be the bearer of
accountability for them.

Kontablo therefore takes the resulting position without hedging: \textbf{the
chain of responsibility is explicit, and it terminates in a human.} We do not
believe --- and the architecture does not assume --- that the \emph{ephemeral
instance} of an AI system, spun up for a single inference and then discarded, can
or should hold legal or professional responsibility for a financial record. The
agent is a proposer and an instrument; whatever responsibility it appears to
exercise is, at every moment, borrowed from and answerable to the human operator
who configured it, authorized its policy, and retains the veto
(Section~\ref{subsec:enforcement_boundary}).

This premise has a consequence that is easy to state and important not to
lose: \textbf{the human role in Kontablo is foundational, not residual.} It is
present in every transaction as the locus of legal accountability, not merely
invoked when the automated tiers fail. We distinguish two analytically
separate human functions, because conflating them produces the false inference
that the human becomes dispensable as automated coverage improves:

\begin{enumerate}
    \item \textbf{Foundational accountability and context governance
    (invariant).} The human professional is the legal bearer of
    responsibility for the record and the guarantor of the \emph{integrity and
    continuity of context}. These are two distinct properties. \emph{Integrity}
    is a point-in-time property: that the jurisdictional rules, effective-date
    regulations, and entity-specific facts the system reasons over are complete
    and correct for the transaction at hand. \emph{Continuity} is a temporal
    property: that this context persists coherently across transactions and
    reporting periods. The distinction is not cosmetic. A stochastic language
    model is \emph{stateless} --- it carries no memory between invocations and
    therefore possesses no temporal continuity of its own. Continuity must
    consequently be supplied by the system rather than the model: it is
    materialized deterministically by the harness (append-only ontology
    versioning, effective-date tracking, and the audit trail described below,
    whose cryptographic sealing is specified as roadmap work) and held to
    account by the human. This function does not diminish as the ontology matures; it applies
    to every transaction regardless of confidence score. Crucially,
    accountability for context integrity and continuity is not the same as
    manual re-verification of every transaction: their \emph{maintenance} is
    largely a systemic, deterministic responsibility of the harness, while
    \emph{accountability} for them remains human. This is the shift from
    point-in-time audit to continuous ``living compliance'' (LSE Business
    Review, 2026).
    \item \textbf{Residual error resolution (reactive).} The human resolves
    the cases that automated tiers cannot --- the boundary population analyzed
    in Section~\ref{subsec:cra_as_error_management}. This load \emph{does}
    decrease as coverage expands.
\end{enumerate}

The architecture below operationalizes both functions. The shrinkage argument
made later in this paper (Section~\ref{subsec:cra_as_error_management}) applies
strictly to function~(2). Function~(1) --- legal accountability and context
governance --- is invariant by construction and is the reason the CRA can
never reduce to zero human involvement, however complete the ontology becomes.

\subsection{The Co-responsibility Architecture (CRA)}
To prevent ``AI-only audit loops''—where automated systems silently hide booking errors or fraud from human supervisors—Kontablo implements a \textbf{Co-responsibility Architecture (CRA)}. This architecture formalizes the governance boundary between human operators and AI models, ensuring that neither party can modify financial records without leaving a verifiable audit trail. The CRA is, in the vocabulary of contemporary \textbf{AI governance}, a domain-specific human-in-the-loop control: it instantiates the accountability, human-oversight, and record-keeping obligations that frameworks such as the NIST AI Risk Management Framework, ISO/IEC~42001, and the EU AI Act place on high-risk AI systems (ISACA, 2024; EU AI Act, 2025), specialized to the accounting ledger where the accountable party is a named professional rather than an organization. Consistent with Section~\ref{subsec:non_delegability}, the AI is positioned throughout as a \emph{proposer and instrumenter} --- it generates informed suggestions, attaches audit and provenance metadata, and improves operational efficiency --- never as the bearer of accountability for the resulting record.

The CRA operates as a multi-stage control loop:
\begin{enumerate}
    \item \textbf{AI Proposes:} The semantic mapping engine analyzes incoming transactions and proposes a Level 3 target node ($V_{\text{universal}}$), attaching a confidence score ($\theta$) and a natural language justification.
    \item \textbf{Human Disposes:} A human accountant or domain expert reviews the proposal. The human has the authority to approve the proposed mapping or override it with a different Level 3 classification.
    \item \textbf{Rule Engine Evaluates:} If the human overrides the proposal, the deterministic rule engine verifies whether the new mapping violates any core ontology constraints (e.g., nature consistency, statement boundaries, or liquidity limits). If a violation is detected, the engine does not block the transaction—which would halt business operations—but instead forces the human to enter a justification.
\end{enumerate}

% Figure: Co-responsibility Architecture (CRA) feedback loop
% Source: originally in sections/governance.tex
% Depicts: AI Proposes → Human Disposes → Rule Engine validates →
%          Inconsistency Flagged + Audit Ledger (immutable).
%
% Design system: local/source=orange!20, Kontablo/universal=blue!15,
%   IFRS/output/approved=green!15, decision=yellow!25, database=gray!20,
%   flag=red!15. Arrows >=Latex, thick. rounded corners=4pt.
\begin{figure}[H]
\centering
\fitwidth{%
\begin{tikzpicture}[
    >=Latex,
    node distance=2cm,
    procnode/.style={draw=blue!60!black, rectangle, rounded corners=4pt,
        fill=blue!10, text width=3.8cm, align=center,
        minimum height=1.2cm, font=\small},
    decnode/.style={draw=yellow!70!black, diamond, aspect=2,
        fill=yellow!20, text width=2.8cm, align=center, font=\small},
    flagnode/.style={draw=red!60!black, rectangle, rounded corners=4pt,
        fill=red!10, text width=3.8cm, align=center,
        minimum height=1.2cm, font=\small},
    dbnode/.style={draw=gray!60, cylinder, shape border rotate=90,
        fill=gray!20, text width=3.0cm, align=center,
        minimum height=1.5cm, font=\small},
    edgelabel/.style={font=\scriptsize, fill=white, inner sep=1pt}
]

% Nodes
\node[procnode] (ai_propose) {%
    \textbf{1.~AI Proposes}\\[2pt]
    Generates mapping proposal and confidence score};
\node[procnode, right=1.8cm of ai_propose] (human_override) {%
    \textbf{2.~Human Disposes}\\[2pt]
    Accountant reviews or overrides mapping};
\node[decnode, below=1.2cm of human_override] (rule_check) {%
    \textbf{3.~Rule Engine}\\Validates deterministic boundaries};
\node[flagnode, left=2.6cm of rule_check] (flag_record) {%
    \textbf{4.~Inconsistency Flagged}\\[2pt]
    Set \texttt{inconsistency\_flag}~= True; write audit note};
\node[dbnode, below=1.2cm of flag_record] (db) {%
    \textbf{Audit Ledger}\\[2pt]
    Persist signed transactional metadata};

% Paths
\draw[->, thick] (ai_propose) -- (human_override);
\draw[->, thick] (human_override) -- (rule_check);
\draw[->, thick] (rule_check) -- node[edgelabel, above] {Violation} (flag_record);
\draw[->, thick] (rule_check.south) -- ++(0,-0.7) -|
    node[edgelabel, near start, right] {Valid} (db.east);
\draw[->, thick] (flag_record) -- (db);

\end{tikzpicture}%
}
\caption{The Co-responsibility Architecture (CRA) feedback loop. An
AI-proposed mapping is reviewed by a human accountant. If the human
override violates deterministic ontology rules (e.g., changing cash
nature), the system forces entry of an explanation, flags the
transaction, and writes it to the immutable ledger for future auditor
review.}
\label{fig:cra_loop}
\end{figure}


\subsection{Database Persistence and Flagging Mechanics}
The outputs of the CRA are persisted directly within the database schema of the ERP integration (e.g., as custom fields in Frappe/ERPNext's \texttt{Journal Entry Account} DocType or SAP B1's user-defined fields). 

Two key metadata fields are appended to every mapped transaction record:
\begin{itemize}
    \item \texttt{inconsistency\_flag}: A boolean field that is set to \texttt{true} if the mapping violates a deterministic rule.
    \item {\sloppy\texttt{inconsistency\_note}: A text field containing the specific rule violation description (e.g., ``Nature Mismatch: local debit account mapped to credit node \texttt{equity.capital}'') concatenated with the human operator's written justification for the override.\par}
\end{itemize}

The specification requires that, once written, these fields be cryptographically hashed and sealed as part of the transaction block, so that any post-facto change to a mapping --- for example, to cover up an irregularity --- produces a hash mismatch in the audit ledger. The v0.x reference implementation provides append-only logging of these fields; cryptographic sealing (hybrid classical and post-quantum signatures over version hashes) is specified as roadmap work and is not yet implemented.

\subsection{Scope and Enforcement Boundary}
\label{subsec:enforcement_boundary}

A recurring question determines whether the CRA is a governance guarantee or a
decoration: \emph{on which transactions does it actually run?} Kontablo is a
\textbf{semantic mapping and validation layer}, not a book of record. The ERP
(ERPNext, Odoo, SAP B1, NetSuite) remains the system of record; Kontablo
resolves each local account to a universal node, runs the deterministic
boundary checks, and writes the resulting mapping and CRA metadata back into
the ERP record (the custom fields of the previous subsection) or returns them
over its API. It does not maintain a shadow ledger the ERP is unaware of.
Three integration modes therefore define the surface the CRA governs:

\begin{enumerate}
    \item \textbf{Connector-mediated.} The connector subscribes to the host
    ERP's write events (e.g.\ Frappe/ERPNext document hooks, SAP B1 UDF
    triggers) so that a posting --- whether typed into the ERP UI, imported
    from a spreadsheet, or pushed by an integration --- is routed through the
    mapping and the boundary checks at write time.
    \item \textbf{API-submitted.} An autonomous agent or upstream service calls
    the mapping API directly (the agent-native path of
    Section~\ref{sec:harness}); the CRA runs inline before the transaction is
    handed to the ERP.
    \item \textbf{Batch / import.} Bulk uploads (CSV exports, migration loads)
    are processed through the import stream
    (Section~\ref{subsec:tree_to_graph}) and validated as a set.
\end{enumerate}

The honest boundary is this: \textbf{the CRA enforces only on transactions that
reach one of these paths.} A journal entry posted directly in the ERP through a
write path the connector does \emph{not} subscribe to is governed only
retroactively --- at the next validation sweep or consolidation run --- not at
the instant of posting. Completeness of enforcement is therefore an integration
property, not an automatic one: it holds exactly to the extent that the
connector intercepts every write path of the host system. We state this
explicitly because conflating ``Kontablo validates mappings'' with ``Kontablo
governs every transaction in the enterprise'' is the ambiguity an implementer
would otherwise discover late. For entries that bypass the live hooks,
Kontablo's guarantee is \emph{detectability on sweep} --- the same
deterministic checks, applied in batch, flag the inconsistency and quarantine
it from the consolidated statement --- not prevention at keystroke.

\paragraph{Two authorization modes, one accountable human.} The CRA does not
require a human to hand-approve every line. Deterministic resolutions (Tier~1
exact-code and Tier~2 rule matches, Section~\ref{sec:deterministic}) post under
\emph{standing ex-ante authorization}: the responsible accountant approves the
deterministic policy --- the code maps and the boundary rules --- once,
configures it, and retains veto and the right to inspect or reverse any posting
it produced. These postings are logged and fully reviewable but are not
individually signed, exactly as an accountant today authorizes a recurring
posting rule without re-signing each instance it fires. Only the residual ---
flagged boundary violations, Tier~3 low-confidence proposals, and explicit
human overrides --- requires \emph{per-transaction disposition}: the human
selects or confirms the mapping and, where a deterministic rule is
contradicted, enters the written justification persisted in
\texttt{inconsistency\_note}. The accountant's ``signature'' is thus two
records, not one: a standing authorization of the deterministic policy, and a
per-exception sign-off on each flagged item. This is what is meant by a
``mandatory human review \emph{pathway} for every proposal'': the pathway and
the veto exist for every transaction, while the act of review is spent where it
has value --- on the exceptions --- rather than rubber-stamping the
deterministic majority. Legal accountability remains undivided and human in
both modes (Section~\ref{subsec:non_delegability}).

\subsection{Anti-Corruption and Embezzlement Mitigation}
In legacy corporate environments, fraud often occurs through the collusion of internal accounting staff who ``rubber-stamp'' misclassified journal entries (e.g., booking personal expenditures as business-related ``Administrative Expenses'' or masking asset withdrawals as ``Intercompany Receivables''). 

By inserting the CRA between the general ledger database and the user interface, Kontablo establishes an automated, independent ``witness.'' Because the deterministic rules cannot be bribed or pressured, any fraudulent mapping override will automatically flag itself in the audit database. When external auditors perform their review, they do not need to sample thousands of transactions randomly; they simply query for all records where \texttt{inconsistency\_flag == true}. This structural transparency dramatically increases the difficulty of internal embezzlement and enforces compliance at the database write-in layer.

\subsection{How does Kontablo relate to sovereign AI?}
\label{subsec:sovereign_ai}
\textbf{Sovereign AI} --- the capacity of a nation or institution to operate and
govern its own AI stack (compute, models, and data) under its own jurisdiction
rather than depending on a foreign provider --- is becoming a load-bearing
constraint for public-sector and regulated finance. Kontablo is
\emph{sovereignty-compatible by construction}, at three levels this paper has
already established and a future revision will develop in full.
\emph{Jurisdictional sovereignty:} the universal layer interoperates with, rather
than replaces, each nation's statutory chart, so adoption requires neither
surrendering a local chart of accounts nor harmonizing sovereign tax codes
(Section~\ref{sec:precedents}). \emph{Data sovereignty:} the per-client
determinism design keeps the agent and its data on client infrastructure
(Section~\ref{subsec:adr011_preview}), so transaction data need not leave the
jurisdiction to be mapped. \emph{Compute and linguistic sovereignty:} because the
deterministic tiers resolve the large majority of volume with no language-model
call at all, a sovereign deployment is not forced onto a foreign frontier model
for its load-bearing decisions --- and the inverse parametric knowledge effect
(Section~\ref{subsec:ipke}) shows that deterministic grounding benefits
non-Anglophone jurisdictions the most. Combined with the non-delegability premise
(Section~\ref{subsec:non_delegability}), under which legal accountability rests on
a professional answerable to national law and never on a borderless AI, these
properties position Kontablo as a candidate substrate for sovereign financial AI.
We surface this point only at a high level here; a fuller treatment of digital
sovereignty --- on-premises and air-gapped deployment profiles, and cross-border
consolidation under conflicting data-residency regimes --- is deferred to future
work (Section~\ref{sec:conclusion}).
